% from main.tex
%%%%%%%%%%  scan idx 528  |  folio 515  %%%%%%%%%%
morphisme structural de $P$, alors $f^*$, défini car $f$ est plat, est surjectif et définit un isomorphisme :
\[
K_\bullet(S)/J\cdot K_\bullet(S)\xrightarrow{\sim}K_\bullet(P).
\]

Soient $E_0=(\mathcal O_S)^n$, $e_1,\ldots,e_n$ la base évidente de $E_0$, $E_0^{(i)}$ le sous-Module de $E_0$ engendré par $e_1,\ldots,e_i$. Le sous-foncteur de $G$ qui à tout $S$-schéma $S'$ associe le sous-groupe des automorphismes de $E_{0S'}$ qui laissent invariants les sous-modules $E_{0S'}^{(i)}$ pour tout $i$ est représentable par un sous-schéma en groupes $B$ de $G$, qui est un sous-groupe de Borel de $G$. Le faisceau quotient $G/B$ (pour la topologie fpqc) est représentable par le schéma des drapeaux de $E_0$, qui sera noté $X_0$. Soient $L_0^{(i)}$ les quotients canoniques de l'image inverse de $E_0$ sur $X_0$, $G_0^{(i)}$ le noyau de $E_{0X_0}\to L_0^{(i)}$, et $F_0^{(i)}=G_0^{(i-1)}/G_0^{(i)}\simeq \Ker(L_0^{(i)}\to L_0^{(i-1)})$.

Par ailleurs, $B$ est produit semi-direct de son radical unipotent $U$ et de son sous-tore maximal $T$ laissant les $e_i$ invariants. Soit $Y_0=G/U$. Alors $Y_0$ est un torseur de base $X_0$ sous le groupe $T_{X_0}$. Si $S'$ est un $S$-schéma, la donnée d'un $S'$-point de $Y_0$ équivaut à la donnée, sur un $S''$ fidèlement plat quasi-compact sur $S'$, d'un automorphisme de $E_{0S''}$ modulo les unipotents de $U_{S''}$, c'est-à-dire d'un drapeau de $E_{0S''}$, les sous-modules du drapeau étant libres, et pour tout $i$, d'une classe d'éléments $f_i^\alpha$ libres sur $\mathcal O_{S''}$ tels que $f_i^\alpha-f_i^\beta\in G_{0S''}^{(i+1)}$, soit encore à la donnée d'un drapeau de $E_{0S''}$ et pour tout $i$ d'une base de $F_{0S''}^{(i)}$, cette donnée étant soumise à la condition de recollement sur $S''\times_{S'}S''$; par descente fpqc, cela revient à la même donnée sur $S'$. On voit alors immédiatement que $Y_0$ est le $X_0$-schéma affine défini par l'algèbre :
\[
A_0=\bigoplus_{k_i\in\mathbb Z}(F_0^{(1)})^{\otimes k_1}\otimes\cdots\otimes(F_0^{(n)})^{\otimes k_n}.
\]
